\documentclass[portuguese]{sbcreviews-2025}

\usepackage[misc,geometry]{ifsym}
\usepackage{aas_macros}
\usepackage[bottom]{footmisc}
\usepackage{tabularray}
\usepackage{url}
\usepackage{cuted}
\usepackage{capt-of}

\UseTblrLibrary{booktabs}

\setcitestyle{square}

\definecolor{unilogo}{rgb}{0.16, 0.26, 0.58}
\definecolor{maillogo}{rgb}{0.58, 0.16, 0.26}
\definecolor{darkblue}{rgb}{0.0,0.0,0.0}
\hypersetup{
    colorlinks,
    breaklinks,
    linkcolor=darkblue,
    urlcolor=darkblue,
    anchorcolor=darkblue,
    citecolor=darkblue
}

\jid{ROCS}
\jtitle{Artigo de Projeto}
\issn{2966-3938}
\doi{}
\copyrightstatement{}
\jyear{2026}

\category{Projeto de Software}
\title[Portal Web da Bamak]{Portal Web Institucional-Comercial para a Bamak: Fundamentação, Proposta, Arquitetura e Cronograma}

\author[Bauer]{
\affil{\textbf{Bredley Bauer}~\textcolor{blue}{\faEnvelope}~~[Centro Universitário Católica de Santa Catarina~|~\href{mailto:bredley.bauer@catolicasc.edu.br}{\textit{bredley.bauer@catolicasc.edu.br}}]}
}

\setlength{\textfloatsep}{0.8em}
\setlength{\floatsep}{0.7em}
\setlength{\intextsep}{0.7em}
\setlength{\emergencystretch}{1.5em}

\begin{document}

\begin{frontmatter}

\maketitle

\begin{mail}
Curso de Engenharia de Software, Centro Universitário Católica de Santa Catarina, Jaraguá do Sul, SC, Brasil.
\end{mail}

\begin{abstract-pt}
Este artigo apresenta a consolidação do projeto de um portal web institucional-comercial para a Bamak Equipamentos LTDA. O trabalho parte do diagnóstico do portal atual, que apresenta a empresa de forma insuficiente, organiza mal informações institucionais e comerciais e oferece pouco apoio ao contato com parceiros e clientes. A partir desse problema, o estudo reúne benchmarking com artigos científicos e referências web do setor, definição da solução proposta, direcionamento arquitetural, resultados esperados e cronograma de desenvolvimento. A solução organiza a presença digital da empresa em torno de apresentação institucional, segmentos, soluções, catálogo de produtos, FAQ sobre orçamento e pedidos, contato, notícias, eventos e um módulo administrativo autenticado para gestão de conteúdo.
\vspace{0.8cm}
\end{abstract-pt}

\end{frontmatter}

\section{Introdução}

Empresas que dependem de credibilidade institucional e de contato frequente com clientes e parceiros precisam apresentar suas informações de forma clara, organizada e acessível no ambiente digital. Quando isso não ocorre, o site deixa de apoiar o posicionamento da empresa e passa a funcionar apenas como presença residual.

A Bamak Equipamentos LTDA atua no fornecimento de materiais, peças e equipamentos para outras empresas, com destaque para soluções voltadas à agroindústria. O portal atual, porém, não comunica com clareza a atuação da empresa, não organiza adequadamente segmentos, soluções e produtos e oferece pouco apoio ao contato comercial.

O problema central do projeto consiste em definir de que forma um portal web institucional-comercial pode reorganizar as informações da Bamak, qualificar sua apresentação institucional e facilitar o contato com potenciais parceiros e clientes. Parte-se da hipótese de que um portal com estrutura informacional mais clara, catálogo de produtos, FAQ sobre orçamento e pedidos e um módulo administrativo autenticado para atualização contínua de conteúdo responde de forma mais adequada à realidade da empresa.

\section{Fundamentação e Benchmarking}

\subsection{Critérios de comparação}

O benchmarking foi organizado para responder duas questões: quais critérios já aparecem nas referências analisadas e quais permanecem como oportunidade de diferenciação no caso da Bamak. Para isso, foram definidos dez critérios diretamente ligados ao problema do projeto: apresentação institucional clara, organização de segmentos e soluções, navegação orientada ao usuário, contato comercial visível, orientação sobre orçamento e pedidos, catálogo de produtos integrado, atualização de conteúdo por painel administrativo, adequação ao contexto B2B, fluxo entre descoberta, entendimento e contato e aderência ao recorte agroindustrial.

Os artigos científicos foram lidos em chave conceitual, avaliando quais critérios são discutidos ou sustentados teoricamente. As referências web foram lidas em chave funcional, observando o que aparece implementado ou sinalizado nas soluções analisadas.

\subsection{Matriz comparativa}

\begin{strip}
\vspace{-0.35em}
\centering
\scriptsize
\captionof{table}{Matriz comparativa entre artigos científicos, referências web e a proposta da Bamak.}
\vspace{0.3em}

{\footnotesize
\parbox{0.98\textwidth}{\centering
\textbf{Referências da matriz:}
Art. A = Chakraborty et al. (2005);
Art. B = Chiou et al. (2010);
Art. C = Ramos et al. (2019);
Web A = site atual da Bamak;
Web B = Texas Equipamentos;
Web C = OBR Equipamentos Industriais.
}
}

\vspace{0.35em}

\begin{tblr}{
    width=\textwidth,
    colspec={Q[l,wd=4.35cm] Q[c,wd=0.95cm] Q[c,wd=0.95cm] Q[c,wd=0.95cm] Q[c,wd=0.95cm] Q[c,wd=0.95cm] Q[c,wd=0.95cm] Q[c,wd=1.05cm]},
    rows={valign=m},
    row{1}={bg=gray!12,font=\bfseries},
    rowsep=3pt,
    hlines,
    vlines
}
Critério analisado & Art. A & Art. B & Art. C & Web A & Web B & Web C & Proposta \\
Apresentação institucional clara & Parcial & Parcial & Parcial & Não & Sim & Sim & Sim \\
Organização de segmentos e soluções & Parcial & Parcial & Não & Não & Sim & Parcial & Sim \\
Navegação orientada ao usuário & Parcial & Sim & Sim & Não & Parcial & Sim & Sim \\
Contato comercial visível & Não & Não & Não & Não & Sim & Sim & Sim \\
Orientação sobre orçamento e pedidos & Não & Não & Não & Não & Não & Parcial & Sim \\
Catálogo de produtos integrado & Não & Não & Não & Parcial & Parcial & Parcial & Sim \\
Atualização de conteúdo por painel administrativo & Não & Não & Não & Não & Não & Não & Sim \\
Adequação ao contexto B2B & Sim & Parcial & Parcial & Parcial & Sim & Sim & Sim \\
Fluxo entre descoberta, entendimento e contato & Parcial & Parcial & Sim & Não & Parcial & Parcial & Sim \\
Aderência ao recorte agroindustrial & Não & Não & Não & Parcial & Não & Não & Sim \\
\end{tblr}

\vspace{0.25em}
{\footnotesize \textbf{Legenda:} Sim = critério atendido; Parcial = critério abordado, mas não resolvido completamente; Não = critério ausente ou não identificado no escopo.}
\vspace{-0.2em}
\end{strip}

\vspace{0.25em}

\subsection{Leitura crítica da matriz}

Os artigos selecionados sustentam critérios relevantes para o projeto, mas não oferecem uma resposta integrada ao problema da Bamak. Chakraborty, Srivastava e Warren \cite{chakraborty2005} reforçam a importância da clareza das informações em websites corporativos B2B, aspecto diretamente aplicável à página inicial, à apresentação institucional e à exposição das soluções da empresa. Chiou, Lin e Perng \cite{chiou2010} ajudam a justificar a reorganização do conteúdo e da estrutura de navegação, sobretudo na hierarquia do menu e na separação entre páginas. Ramos, Rita e Moro \cite{ramos2019} sustentam a centralidade da usabilidade, o que, neste caso, significa reduzir esforço para que o visitante encontre segmentos, soluções, produtos, orientação sobre pedidos e contato.

Na comparação com as referências web, o site atual da Bamak \cite{bamaksite} aparece como ponto de partida do problema, pois atende de forma insuficiente a critérios centrais da matriz. Texas Equipamentos \cite{texassite} e OBR Equipamentos Industriais \cite{obrsite} resolvem melhor aspectos específicos, como apresentação institucional, navegação e contato, mas não articulam, no escopo observado, uma combinação completa entre organização de segmentos, soluções e produtos, orientação sobre orçamento e pedidos, atualização contínua de conteúdo por painel administrativo e recorte agroindustrial.

A proposta da Bamak se diferencia pela combinação desses elementos em uma mesma estrutura. O diferencial não está em um recurso isolado, mas na articulação entre apresentação institucional orientada ao público corporativo, organização explícita de segmentos, soluções e produtos, FAQ como ponte entre interesse e contato, agenda pública de eventos e feiras e gestão de conteúdo em módulo administrativo autenticado.

\section{Proposta do Portal}

A solução proposta para a Bamak foi organizada em duas frentes complementares: uma área pública e um módulo administrativo autenticado. A área pública concentra a apresentação institucional da empresa e o percurso principal do visitante. Nela, o portal deve reunir página inicial, apresentação da empresa, seções de segmentos e soluções, catálogo de produtos organizado por categoria ou segmento, FAQ sobre orçamento e pedidos, página de contato, notícias e eventos e feiras em destaque.

Essa composição atende a um usuário que chega ao site sem conhecer a empresa e precisa entender rapidamente quem é a Bamak, em que frentes atua, que tipo de solução oferece e como avançar para o contato comercial. O catálogo de produtos deixa de funcionar como vitrine isolada e passa a se integrar à lógica de segmentos e soluções. A FAQ, por sua vez, não substitui o contato, mas funciona como etapa de orientação antes dele.

O módulo administrativo foi definido como parte formal do sistema. Seu papel é sustentar a atualização contínua do portal e evitar que o conteúdo volte a ficar estático. Nessa área, estão previstos login administrativo, painel inicial simples, gestão de notícias, eventos e feiras, FAQ, catálogo de produtos e conteúdos institucionais centrais. Também está prevista a visualização das mensagens recebidas pelo formulário de contato.

Com base nesse escopo, o projeto passa a operar em quatro fluxos principais: navegação institucional, contato orientado por FAQ, gestão de notícias e gestão de eventos e catálogo. Esses fluxos estruturam a solução em nível funcional e servem como base para o mapeamento de telas, a prototipação e o desenvolvimento posterior.

\section{Arquitetura e Tecnologias}

A arquitetura da solução foi definida em modelo cliente-servidor, com separação entre interface web, processamento das regras de negócio e persistência de dados. No nível de contexto, o sistema se relaciona com dois perfis de uso: o visitante ou parceiro comercial, que acessa o portal para conhecer a empresa, consultar segmentos, soluções, catálogo, FAQ, notícias e eventos, e o administrador da Bamak, responsável pela atualização dos conteúdos institucionais e comerciais. O sistema também se integra a um serviço externo de e-mail para encaminhamento das mensagens enviadas pelo formulário de contato.

No nível de containers, a solução foi organizada em três partes principais: uma aplicação web frontend, uma API backend e um banco de dados relacional. A aplicação frontend concentra a interface pública do portal e o acesso à área administrativa autenticada. A API backend centraliza autenticação administrativa, regras de negócio e operações de leitura e escrita dos conteúdos gerenciados no sistema. O banco de dados armazena usuários administrativos, notícias, eventos, itens de FAQ, produtos do catálogo e mensagens de contato.

Para o frontend, foi definido o uso de Next.js com TypeScript. Essa combinação é adequada ao projeto por sustentar uma interface web contemporânea, componentizada e compatível com a necessidade de evolução futura do portal. Na camada visual, foi adotado Tailwind CSS, com componentes personalizados estruturados a partir de shadcn/ui. Essa decisão mantém controle sobre a identidade da interface e evita dependência de um conjunto visual fechado.

No backend, foi definido o uso de NestJS com TypeScript. A escolha é compatível com a necessidade de estruturar autenticação, gerenciamento de conteúdo e integração entre os módulos públicos e administrativos em uma API organizada por responsabilidades. Para persistência, foi definido PostgreSQL como banco de dados relacional e Prisma ORM como camada de acesso a dados. Essa combinação atende à modelagem prevista para o projeto, que envolve entidades bem definidas e relacionamentos entre conteúdos institucionais, produtos, mensagens e usuários administrativos.

A autenticação da área administrativa será baseada em JWT, restringindo o acesso ao módulo de gestão de conteúdo. Também foi definida a adoção de Docker como estratégia de padronização de ambiente para desenvolvimento e implantação futura. Em conformidade com as diretrizes da linha Web Apps, a proposta mantém backend próprio, banco de dados relacional e estrutura explícita de execução, sem depender de plataformas no-code ou soluções que ocultem a arquitetura da aplicação.

Assim, a stack final do projeto ficou organizada da seguinte forma: Next.js, TypeScript, Tailwind CSS e componentes personalizados baseados em shadcn/ui no frontend; NestJS, TypeScript e Prisma no backend; PostgreSQL na persistência de dados; JWT na autenticação administrativa; e Docker na padronização de ambiente. Os diagramas C1 e C2 elaborados para o projeto formalizam essa definição ao representar, respectivamente, o contexto de uso do sistema e sua organização interna em containers.

\section{Resultados Esperados, Discussão e Cronograma}

O primeiro resultado esperado é a reorganização efetiva da presença digital da Bamak. Isso significa tornar mais clara a apresentação institucional da empresa, explicitar seus segmentos e soluções, estruturar melhor a exposição de produtos e reduzir a distância entre a descoberta inicial do portal e o contato comercial.

O segundo resultado esperado está na manutenção do conteúdo. A proposta inclui um módulo administrativo autenticado para notícias, eventos, FAQ, catálogo e conteúdos institucionais centrais. Com isso, o site deixa de depender exclusivamente de atualização estática e passa a operar com um nível mínimo de continuidade editorial.

O terceiro resultado esperado está na coerência entre escopo e necessidade real. O projeto não foi direcionado para orçamento automatizado, e-commerce, CRM ou área de cliente. Essa delimitação é intencional. O problema identificado na empresa não exige, neste momento, um sistema transacional complexo; exige um portal que organize melhor a comunicação institucional e comercial e sustente atualização contínua de conteúdo.

Em termos acadêmicos, a proposta também produz uma base consistente para a próxima etapa do trabalho. A definição funcional, a arquitetura, os fluxos do sistema e a prototipação prevista permitem que o desenvolvimento futuro comece com menos ambiguidade e maior rastreabilidade entre benchmarking, decisão de projeto e implementação.

O cronograma de desenvolvimento foi organizado em quatro etapas principais. A primeira concentra arquitetura, requisitos, arquitetura da informação e prototipação em baixa fidelidade. A segunda concentra guia de interface, prototipação em média fidelidade e consolidação do artigo. A terceira contempla validação com a empresa e ajustes finais de conteúdo. A quarta reúne finalização do relatório e preparação da apresentação. Essa divisão permite encadear planejamento, validação e fechamento documental sem romper a coerência da proposta.

\section{Considerações Finais}

O artigo consolidou a proposta de um portal web institucional-comercial para a Bamak a partir de três frentes: diagnóstico do problema, benchmarking e definição da solução. O problema identificado não está na ausência de presença digital, mas na baixa capacidade do portal atual de representar a empresa com clareza e apoiar contato comercial de forma consistente.

O benchmarking mostrou que referências acadêmicas e referências web do setor ajudam a sustentar critérios importantes para a solução, mas também evidenciam lacunas que justificam a proposta adotada. A resposta definida para a Bamak combina apresentação institucional, organização de segmentos, soluções e produtos, orientação sobre orçamento e pedidos, agenda pública de eventos e gestão administrativa de conteúdo.

A definição arquitetural complementa esse percurso ao transformar a proposta em uma solução tecnicamente delimitada. Com isso, o projeto encerra esta etapa com um escopo funcional claro, uma arquitetura coerente com as exigências da universidade e uma base adequada para a prototipação e o desenvolvimento do próximo semestre.

\begin{thebibliography}{9}
\small
\setlength{\itemsep}{0.25em}
\setlength{\parskip}{0pt}

\bibitem{chakraborty2005}
Chakraborty, G.; Srivastava, P.; Warren, D. L. (2005).
\textit{Understanding corporate B2B web sites' effectiveness from North American and European perspective}.
Industrial Marketing Management, 34(5), 420--429.

\bibitem{chiou2010}
Chiou, W.-C.; Lin, C.-C.; Perng, C. (2010).
\textit{A strategic framework for website evaluation based on a review of the literature from 1995--2006}.
Information \& Management, 47(5--6), 282--290.

\bibitem{ramos2019}
Ramos, R. F.; Rita, P.; Moro, S. (2019).
\textit{From institutional websites to social media and mobile applications: A usability perspective}.
European Research on Management and Business Economics, 25(3), 138--143.

\bibitem{bamaksite}
Bamak Equipamentos LTDA.
\textit{Site oficial da Bamak Equipamentos LTDA}.
Disponível em: \url{https://bamak.com.br/}. Acesso em: 14 abr. 2026.

\bibitem{texassite}
Texas Equipamentos.
\textit{Texas Equipamentos -- Institucional}.
Disponível em: \url{https://www.texas.com.br/institucional}. Acesso em: 14 abr. 2026.

\bibitem{obrsite}
OBR Equipamentos Industriais.
\textit{OBR Equipamentos Industriais -- Página inicial}.
Disponível em: \url{https://www.obr.com.br/}. Acesso em: 14 abr. 2026.

\end{thebibliography}

\end{document}